\documentclass[11pt,a4paper]{article}
\usepackage[margin=2.2cm]{geometry}
\usepackage{lmodern}
\usepackage[T1]{fontenc}
\usepackage{booktabs,array,longtable,xcolor,enumitem}
\usepackage[hidelinks]{hyperref}
\newcommand{\TBD}[1]{\textcolor{red}{[TBD\ifx\relax#1\relax\else: #1\fi]}}
\setlist{nosep,leftmargin=*}
\setlength{\parskip}{0.4em}\setlength{\parindent}{0pt}
\newcommand{\concern}[2]{\subsection*{#1\quad #2}}
\title{Response to the review of\\ ``Baseline Configuration and Random-Sampling Controls in Metaheuristic Comparisons for Wind Farm Layout Optimization''}
\author{Prince Solanki, Pulkit Dwivedi, Vanita Garg, Vinay Shukla}
\date{Revision 3, October 2026}
\begin{document}
\maketitle

We thank the reviewer for the careful report (items R1--R10) and the detailed list of remaining work (Sections A--E).
Below we answer every item. Section, table and equation numbers refer to the revised manuscript (M) and the
revised supplementary material (S). For each item we state whether it is \emph{resolved} by new analyses, new
experiments or text changes, or \emph{disclosed} (a limitation stated in the paper but not removed).

\textbf{A correction of premise that affects several items.} The previous revision stated that the 36,190 original run
records and fifteen core modules were missing from the archive, so that the historical checks and several requested
re-analyses could not be rerun. This was incorrect: the complete archive exists. For this revision we re-executed the
entire analysis pipeline from it in the environment of the runs (CPython 3.11.15, NumPy 2.4.6, SciPy 1.17.1, pandas
3.0.6). All 91 regenerated tables, number files, summaries and figures are identical to those of the manuscript apart
from time stamps. All automated checks pass (62 main-result, 57 statistical, 22 diagnostic, 41 direction-resolution,
15 coefficient-sweep and 57 theory checks). The analysis of the first-revision experiments reproduces its 4,670
values exactly, and reruns of stored runs with the archived drivers are bit-identical except for elapsed time. This
made every item of Section~B and item~C4 feasible. We also added 5,100 new optimization runs (C1, C2) and the
full-precision reruns of C4. All statements about missing material have been removed (M 6.4, Data availability,
S9, S12).

\section*{Reviewer items R1--R10}

\concern{R1}{Unfinished declarations.}
\emph{Resolved in the text; author confirmation pending.} The CRediT roles, competing-interest, funding and
data-availability statements are complete. The declaration on generative AI now names all tools used, in Elsevier's
template wording. The only remaining \texttt{[TBD]} markers are facts that only the authors can supply: the
deposit DOI, confirmation of the AI tool list, and the cover-letter confirmation of originality and approval.

\concern{R2}{Public availability.}
\emph{Resolved up to the public deposit.} A versioned, deposit-ready reproducibility package (release 1.0.0) contains
all run records (36,190 original, 11,640 first-revision and 5,100 new runs, plus the full-precision reruns), all
models, optimizers, drivers, analysis and check scripts, an environment specification (requirements, Dockerfile) and
SHA-256 checksums. Its README maps every table, figure and number to the command that regenerates it. We verified it
from a clean unpack: checksums pass, and the full regeneration reproduces all outputs. The Data availability statement
now describes the package and states that the public DOI will be added before publication. Until then, the package
is available to the editor and reviewers. Changes: M 6.4, Data availability, S9 (Reproducibility Map), S12.

\concern{R3}{Stronger baselines outside the principal tables.}
\emph{Resolved.} The site comparisons were recomputed in single method pools:
\begin{itemize}
\item \emph{Horns Rev~1.} The ten methods plus the GA (Table~10 now names its pool and gives ranks; Table~S71).
PSO-VNS is best at 6,030 evaluations (GA second, $p_{\rm Holm}=0.53$). The GA is best at 30,030 and beats all ten
other methods ($p_{\rm Holm}\le1.1\times10^{-4}$).
\item \emph{IEA37.} The ten methods, the GA and both exact-gradient methods (Tables~11--12, S72). An exact-gradient
method ranks first in every scenario and budget. Among the gradient-free methods, PSO-VNS leads with 16 turbines and
the GA with 36.
\end{itemize}
Rank values from different pools are no longer mixed (M 9.3, 9.5).

\concern{R4}{Conditional ranking versus reliability.}
\emph{Resolved.}
\begin{itemize}
\item \emph{All-run endpoint.} The paired endpoint (feasible beats infeasible, then the objective) is now reported for
all 68 benchmark cases, with case-bootstrap and two-stage intervals and Holm-corrected sign tests (M 9.6, Table~S66).
\begin{itemize}
\item It agrees with the conditional ranks except that SSA and MS-SLSQP swap places.
\item Against PSO, PSO-VNS wins 950, ties 348 and loses 742 of 2,040 seed pairs (score 0.551 [0.516, 0.587]): it is
better more often than not, by amounts within the margin.
\end{itemize}
\item \emph{Imputed cases.} Dropping the imputed ``maximal-difference'' cases, or keeping zero differences (Pratt),
changes no verdict (M 6.2, Table~S69).
\end{itemize}
Feasibility rates, conditional means and the all-run outcome are reported as separate outcomes throughout.

\concern{R5}{Penalty and repair confounding.}
\emph{Resolved by a new experiment; the conclusion is now stated as partly dependent on constraint handling.}
\begin{itemize}
\item \emph{Design.} Pre-registered, 2,700 runs: PSO-VNS, PSO, GA, SSA-VNS and DE, six cases, new seeds 31--60.
Three constraint handlings: (i) the published penalty with box clipping; (ii) Deb's feasibility rules on
dimensionless violations with box clipping; (iii) the same rules with radial projection (M 10.4, Tables~S77--S78).
\item \emph{Validation.} Variant~(i) reproduces 60 of 60 stored runs bit for bit.
\item \emph{With projection.} The rank order is unchanged and every method improves by 0.46--0.80~pp. PSO-VNS keeps
its paired advantage over GA, SSA-VNS and DE, but not over PSO.
\item \emph{With box clipping and normalized violations.} A turbine clipped to the square becomes cheap, and PSO's
feasibility falls from 98.9\% to 73.3\%. PSO-VNS is then no longer distinguishable from GA and SSA-VNS, and the GA
has the best case-level rank.
\end{itemize}
The abstract, contributions, Discussion (11), Threats (12), the checklist (Table~15, new row) and the highlights
now say that the conclusions concern the implemented penalty and repair, and how they depend on it.

\concern{R6}{Proposition~1(c).}
\emph{Resolved in the previous revision and re-verified.} The theory check (57 conditions, including the corrected
statement $V_\infty=0$ iff $c_1c_2(p-g)=0$) passes on the current text.

\concern{R7}{Inference scope, weighting and multiplicity.}
\emph{Resolved.}
\begin{itemize}
\item \emph{Calibration (B3).} The bootstrap TOST $p$ values and the 90\%-interval rule agree for all 18 pairs at every
level. A null simulation at the margin gives empirical sizes of 4.4--7.4\% for PSO-VNS vs.\ PSO. For the case-level
non-superiority test of SSA-VNS vs.\ RSD-VNS it gives 15.0\%, because of one outlying case. Calibrated against the
null, both primary claims still hold ($p=0.015$ and $0.007$; M 6.2, Table~S67, S6.4).
\item \emph{Synchronized seeds (B4).} Seed~$k$ gives identical initial populations across cases. Joint resampling of
one seed vector for all cases leaves PSO-VNS vs.\ PSO unchanged. The seed-level equivalence of SSA-VNS vs.\ RSD-VNS is
then lost, but its non-superiority holds; Table~7 now says so (Table~S68).
\item \emph{Equal-cluster target (B7).} Recomputed from the records, it now has wild-cluster bootstrap intervals; no
verdict changes (Table~S70).
\end{itemize}

\concern{R8}{Re-evaluation versus re-optimization.}
\emph{Resolved by a new experiment.}
\begin{itemize}
\item \emph{Direct 1$^\circ$ optimization (C2).} Pre-registered, 1,200 runs: five methods on eight cases optimized
the 1$^\circ$ objective itself, next to a seed-paired control arm optimized at 15$^\circ$ (1,200 runs; M 10.5,
Table~S79).
\begin{itemize}
\item PSO-VNS ranks first (1.12).
\item Its advantage over PSO grows to $-0.194$~pp (90\% CI $[-0.272,-0.114]$ at the case level), beyond the margin.
\item Optimizing at 1$^\circ$ lowers the 1$^\circ$ wake loss by 0.32--1.12~pp.
\item The leader does not change.
\end{itemize}
\item \emph{GA layouts under other evaluators (C3).} The stored GA Horns Rev~1 layouts were re-evaluated at 1$^\circ$
and with PyWake. GA's lead at 30,030 evaluations holds under both (Table~S73).
\item \emph{Lillgrund at 1$^\circ$ (C3).} The Lillgrund layouts were re-evaluated at 1$^\circ$ as well. PSO-VNS keeps
the highest mean at 30,030, but its lead over VNS and MS-SLSQP is then not significant; the text says so (M 9.4,
Table~S74).
\end{itemize}

\concern{R9}{Charged evaluations versus elapsed time.}
\emph{Resolved.} Elapsed times now come from the original records (Tables~S75--S76, M 9.7).
\begin{itemize}
\item \emph{MS-SLSQP per-evaluation time.} The value we reported (0.53~ms, 2.3$\times$ PSO; median) is the typical
cost. The previously printed mean (1.82~ms) reflects thread contention in one run batch for $N\ge8$: single-threaded
re-timing gives 4--7~s instead of 16--22~s per run.
\item \emph{Feasibility-preserving initialization.} It is not charged to the budget. It takes a median of 4.8~s per run,
more than a 6,030-evaluation PSO-VNS run.
\item \emph{Total compute.} The 36,190 original runs took 101 CPU-hours, run as four parallel processes.
\end{itemize}

\concern{R10}{Attribution, scope and table density.}
\emph{Resolved.} The attribution and scope changes of the previous revision are kept.
\begin{itemize}
\item The principal tables name their method pools and mark the best method.
\item The IEA37 means table is half its former height.
\item Repetition was removed throughout (E1).
\item The bibliographies are in order of first citation, and two references were corrected after a full audit (A6).
\end{itemize}

\section*{Remaining-work report, Sections A--E}
\begin{longtable}{>{\raggedright\arraybackslash}p{0.07\textwidth}>{\raggedright\arraybackslash}p{0.14\textwidth}>{\raggedright\arraybackslash}p{0.71\textwidth}}
\toprule
Item & Status & What was done / what remains \\ \midrule
\endhead
A1 & author & Declarations complete; every author must confirm roles, interests, funding, affiliations. \\
A2 & author & Availability statement now true to the package; public deposit (DOI) before publication. \\
A3 & author & Two primary equivalence questions stated (M 6.2); co-author agreement to be recorded. \\
A4 & author & AI declaration names Claude and ChatGPT; authors confirm the list. \\
A5 & partly & Solanki and Deep (2023) is still Online First (no volume/pages assigned); entry updated accordingly. \\
A6 & done & All 88 main and 23 supplementary references audited; corrections applied; renumbered by first citation. Crossref/publisher pages were not reachable from the audit environment; search-engine records of them were used. \\
A7 & partly & Abstract 239 words ($\le$250), 5 highlights of 64--76 characters ($\le$85), 7 keywords; live Guide for Authors not reachable — authors to confirm. \\
A8 & done/author & Archive and precision items resolved (this revision); author checks remain. \\
B1 & done & All-run endpoint on 68 cases (M 9.6, S66). \\
B2 & done & Horns Rev~1 and IEA37 pools with GA and exact gradients (M 9.3, 9.5; S71, S72). \\
B3 & done & TOST audit and null simulation (M 6.2; S67; S6.4). \\
B4 & done & Synchronized seed resampling (M 6.2, Table~7; S68). \\
B5 & done & Imputation sensitivity (M 6.2; S69). \\
B6 & done & Elapsed times from original records; initialization timing (M 9.7; S75, S76). \\
B7 & done & Wild-cluster bootstrap for the equal-cluster target (S70). \\
C1 & done & Constraint-handling study, 2,700 runs (M 10.4; S77, S78). \\
C2 & done & Direct 1$^\circ$ optimization, 1,200 runs + 1,200-run control (M 10.5; S79). \\
C3 & done & GA (and Lillgrund) layouts at 1$^\circ$ and PyWake (M 9.3, 9.4; S73, S74). \\
C4 & done & All 56,540 stored records audited with the correct site geometry: 50,826 labels confirmed from the rounded coordinates, none contradicted, 5,714 undecidable within the rounding bound; \TBD{C4 rerun: N bit-identical, labels confirmed/changed} (M 9.6; S80, S12.2). \\
D & done/author & Reproducibility package built and verified; public DOI pending (author action). \\
E1 & done & Length pass: Sections 2, 3, 6, 7, 11--13 shortened; the new results are reported in single sentences with supplementary tables; the manuscript stays at 50 pages despite five new analyses. \\
E2 & done & Tables name their pools and mark the best method; no overfull boxes. \\
E3 & done & Supplement has 0 overfull boxes. \\
E4 & done & All PDFs rebuilt from the final sources. \\
E5 & done & Highlights and cover letter updated. \\
E6 & done & This letter. \\
\bottomrule
\end{longtable}

\section*{Items disclosed rather than resolved}
\begin{itemize}
\item The constraint-handling study covers six cases and five methods, and the direct 1$^\circ$ optimization covers
eight cases (M 12).
\item The equivalence margin is post hoc. The cluster-level analyses rest on six designed clusters. All case-level
$p$ values describe this benchmark, not a population of farms (M 6.2, 12).
\item The pool lacks CMA-ES, L-SHADE and tuned DE. The 64-turbine IEA37 scenario and the full Horns Rev~1 and
Lillgrund farms were not run (M 12, 13).
\item The public deposit, the author confirmations and the volume and pages of Solanki and Deep (2023) are author
actions.
\end{itemize}

\end{document}
